\nolabel{section}{Заключение}

В ходе производственной практики была разработана библиотека \texttt{cast-castle}, реализующая метапрограммирование на этапе компиляции для автоматической генерации мапперов с использованием KSP. Библиотека решает актуальную проблему сокращения шаблонного кода при преобразовании моделей данных в чистой архитектуре \cite{refd}. Основные результаты работы:

\begin{itemize}
    \item Разработана аннотация \texttt{@CastCastleMapper}, позволяющая маркировать классы, функции-члены, функции расширения и companion-объекты.
    \item Реализован KSP-процессор, который анализирует структуру исходных и целевых классов, определяет совпадающие и недостающие поля, генерирует вспомогательные функции с добавленными параметрами.
    \item Реализована поддержка вложенных объектов: процессор рекурсивно вызывает мапперы для полей, типы которых также являются моделями.
    \item Реализована поддержка коллекций: для типов \texttt{List}, \texttt{Set} и их мутабельных аналогов генерируется код, преобразующий каждый элемент с помощью соответствующего маппера.
    \item Реализована поддержка функции-расширений: приёмник \texttt{this} используется как исходный объект.
    \item Работа генерации реализована полностью на этапе компиляции без использования рефлексии, что обеспечивает высокую производительность в рантайме и совместимость с ProGuard \cite{refc}.
    \item Библиотека протестирована на примерах из папки \texttt{samples}, включающих различные сценарии использования.
\end{itemize}

Разработанная библиотека может быть использована в реальных Android-проектах для автоматического создания мапперов между слоями, что позволяет разработчикам сосредоточиться на бизнес-логике, а не на рутинном копировании полей \cite{refh}. В дальнейшем планируется добавить поддержку кастомных преобразований типов, более широкую поддержку generic-типов и интеграцию с популярными DI-фреймворками (например, Dagger \cite{refc}, Koin \cite{refg}).
